\documentclass[12pt,letterpaper]{article}
\usepackage[utf8]{inputenc}
\usepackage[spanish,es-noshorthands]{babel}
\usepackage{amsmath,amssymb}
\usepackage{graphicx}
\usepackage{xcolor}
\usepackage{tikz}
\usetikzlibrary{calc,arrows.meta,decorations.markings}
\usepackage[american]{circuitikz}
\usepackage{enumerate}
\usepackage{array}
\usepackage[top=2.5cm, bottom=2.7cm, left=2cm, right=2cm, footskip=40pt]{geometry}
\usepackage{fancyhdr}
\usepackage{lastpage}

\definecolor{ubbblue}{RGB}{0,51,102}

% signo + junto a la placa larga (borne positivo) de la fuente
\newcommand{\bplus}[1]{\node[font=\footnotesize, inner sep=0pt] at (#1) {$+$};}
% malla recorrida en sentido horario: \mallah{x,y}{etiqueta}
\newcommand{\mallah}[2]{%
  \draw[blue, dashed, very thick, postaction={decorate},
        decoration={markings, mark=at position 1 with {\arrow[blue,scale=1.2]{>}}}]
    (#1)+(90:0.5) arc[start angle=90, end angle=-180, radius=0.5];
  \node[blue] at (#1) {#2};}

% ===== PIE DE PAGINA (formato institucional) =====
\pagestyle{fancy}
\fancyhf{}
\renewcommand{\headrulewidth}{0pt}
\renewcommand{\footrulewidth}{0.4pt}
\fancyfoot[L]{\raisebox{-0.5\height}{\includegraphics[height=1.1cm]{logo ubb.png}}}
\fancyfoot[C]{%
    \begin{minipage}[c]{8cm}\centering
        \footnotesize\textcolor{ubbblue}{\textbf{Dr.\ Arturo Fern\'andez P\'erez}}\\
        \scriptsize Departamento de F\'isica -- Universidad del B\'io-B\'io\\
        \scriptsize \textit{Electro\_Asistente\_AI-UBB}\\
        \tiny P\'agina \thepage\ de \pageref{LastPage}
    \end{minipage}}
\fancyfoot[R]{\raisebox{-0.5\height}{\includegraphics[height=1.1cm]{logo dfis.png}}}
\fancypagestyle{plain}{%
    \fancyhf{}
    \renewcommand{\headrulewidth}{0pt}
    \renewcommand{\footrulewidth}{0.4pt}
    \fancyfoot[L]{\raisebox{-0.5\height}{\includegraphics[height=1.1cm]{logo ubb.png}}}
    \fancyfoot[C]{%
        \begin{minipage}[c]{8cm}\centering
            \footnotesize\textcolor{ubbblue}{\textbf{Dr.\ Arturo Fern\'andez P\'erez}}\\
            \scriptsize Departamento de F\'isica -- Universidad del B\'io-B\'io\\
            \scriptsize \textit{Electro\_Asistente\_AI-UBB}\\
            \tiny P\'agina \thepage\ de \pageref{LastPage}
        \end{minipage}}
    \fancyfoot[R]{\raisebox{-0.5\height}{\includegraphics[height=1.1cm]{logo dfis.png}}}
}

\begin{document}

%========================================================
% ENCABEZADO INSTITUCIONAL
%========================================================
\begin{center}
    \includegraphics[height=2cm]{logo ubb.png}\\[0.3cm]
    {\large\textcolor{ubbblue}{\textbf{UNIVERSIDAD DEL B\'IO-B\'IO}}}\\[0.1cm]
    {\normalsize Facultad de Ciencias -- Departamento de F\'isica}
\end{center}

\vspace{0.3cm}\hrule height 1pt\vspace{0.5cm}

{\centering\Large\textbf{Pauta Tarea 3 -- F\'isica II (2026-2)}\par}
{\centering\normalsize\textit{Condensadores, Leyes de Kirchhoff y Fuerza Magn\'etica}\par}

\vspace{0.6cm}

%========================================================
% PROBLEMA 1  (microfaradios, 2 mallas)
%========================================================
\section*{Problema 1 -- Circuito mixto de seis condensadores}

\subsection*{Enunciado}

En el circuito de la Fig.\ 1, una fuente $V = 24{,}0\ [\text{V}]$ alimenta seis condensadores
con capacitancias $C_1 = 12{,}0\ [\mu\text{F}]$, $C_2 = 6{,}0\ [\mu\text{F}]$,
$C_3 = 4{,}5\ [\mu\text{F}]$, $C_4 = 3{,}0\ [\mu\text{F}]$, $C_5 = 6{,}0\ [\mu\text{F}]$ y
$C_6 = 4{,}0\ [\mu\text{F}]$. Determine:

\begin{enumerate}[(a)]
    \item La capacitancia equivalente del circuito.
    \item La carga y la diferencia de potencial de cada condensador.
\end{enumerate}

\begin{center}
\begin{circuitikz}[every node/.style={font=\small}]
  % fuente (+ arriba) entre a y f
  \draw (0,4) to[battery1, l_=$V$] (0,0);
  \bplus{0.32,2.25}
  % riel superior: a -- C1 -- b -- C2 -- c
  \draw (0,4) to[C, l=$C_1$] (4,4) to[C, l=$C_2$] (8,4);
  % rama central b -- C3 -- d  y  lado derecho c -- C4 -- e
  \draw (4,4) to[C, l=$C_3$] (4,0);
  \draw (8,4) to[C, l=$C_4$] (8,0);
  % riel inferior: f -- C6 -- d -- C5 -- e
  \draw (0,0) to[C, l_=$C_6$] (4,0) to[C, l_=$C_5$] (8,0);
  % nodos
  \node[circ] at (4,4) {}; \node[circ] at (4,0) {};
  \node[above left]  at (0,4) {$a$};
  \node[above right] at (4,4) {$b$};
  \node[above right] at (8,4) {$c$};
  \node[below left]  at (0,0) {$f$};
  \node[below right] at (4,0) {$d$};
  \node[below right] at (8,0) {$e$};
\end{circuitikz}\\[2pt]
{\footnotesize Fig.\ 1. Circuito mixto de seis condensadores (dos mallas).}
\end{center}

\subsection*{Datos}

\[
    C_1 = 12{,}0\ [\mu\text{F}],\quad C_2 = 6{,}0\ [\mu\text{F}],\quad C_3 = 4{,}5\ [\mu\text{F}],\quad
    C_4 = 3{,}0\ [\mu\text{F}],\quad C_5 = 6{,}0\ [\mu\text{F}],
\]
\[
    C_6 = 4{,}0\ [\mu\text{F}],\qquad V = 24{,}0\ [\text{V}]
\]

\subsection*{An\'alisis del circuito}

Los puntos $b$ y $d$ son \textbf{nodos}: en cada uno se unen tres conductores. En cambio, $c$ y
$e$ son simples \textbf{esquinas}, donde s\'olo se conectan dos condensadores. Por lo tanto:
\begin{itemize}\setlength\itemsep{1pt}
    \item $C_2$, $C_4$ y $C_5$ forman un \'unico camino en \textbf{serie} entre $b$ y $d$
          ($b\to c\to e\to d$).
    \item Ese camino y $C_3$ conectan los mismos nodos $b$ y $d$: est\'an en \textbf{paralelo}.
    \item $C_1$ ($a$--$b$), el grupo $b$--$d$ y $C_6$ ($d$--$f$) quedan en \textbf{serie} con la
          fuente.
\end{itemize}
\[
    C_{eq}:\qquad C_1 \;-\; \underbrace{\bigl[\,C_3 \parallel (C_2 - C_4 - C_5)\,\bigr]}_{\text{entre } b \text{ y } d} \;-\; C_6
    \qquad(\text{``}-\text{''}:\ \text{serie})
\]

\subsection*{Parte (a) -- Capacitancia equivalente}

\textbf{Paso 1.} $C_2$, $C_4$ y $C_5$ en serie forman $C_a$:
\begin{align*}
    \frac{1}{C_a} &= \frac{1}{C_2} + \frac{1}{C_4} + \frac{1}{C_5}
                   = \frac{1}{6{,}0} + \frac{1}{3{,}0} + \frac{1}{6{,}0}\ [\mu\text{F}^{-1}]\\[4pt]
                  &= \frac{4}{6}\ [\mu\text{F}^{-1}]
                  \;\;\Rightarrow\;\; C_a = 1{,}5\ [\mu\text{F}]
\end{align*}

\textbf{Paso 2.} $C_3$ y $C_a$ en paralelo forman el grupo $C_A$ (entre $b$ y $d$):
\[
    C_A = C_3 + C_a = 4{,}5 + 1{,}5 = 6{,}0\ [\mu\text{F}]
\]

\textbf{Paso 3.} $C_1$, $C_A$ y $C_6$ en serie:
\begin{align*}
    \frac{1}{C_{eq}} &= \frac{1}{C_1} + \frac{1}{C_A} + \frac{1}{C_6}
                      = \frac{1}{12{,}0} + \frac{1}{6{,}0} + \frac{1}{4{,}0}\ [\mu\text{F}^{-1}]\\[4pt]
                     &= \frac{1 + 2 + 3}{12} = \frac{1}{2}\ [\mu\text{F}^{-1}]
\end{align*}
\[
    \boxed{C_{eq} = 2{,}0\ [\mu\text{F}]}
\]

\begin{center}
\begin{circuitikz}[every node/.style={font=\small}]
  \draw (0,2.4) to[battery1, l_=$V$] (0,0);
  \bplus{0.32,1.45}
  \draw (0,2.4) to[C, l=$C_1$] (3,2.4) to[C, l=$C_A$] (6,2.4) -- (6,0)
        to[C, l=$C_6$] (0,0);
\end{circuitikz}\\[2pt]
{\footnotesize Circuito equivalente en serie: $C_1$, $C_A = C_3 \parallel C_a$ y $C_6$.}
\end{center}

\subsection*{Parte (b) -- Carga y diferencia de potencial de cada condensador}

\textbf{Carga total.} Los tres bloques en serie ($C_1$, $C_A$, $C_6$) llevan la misma carga, igual a
la que entrega la fuente:
\[
    Q = C_{eq}\,V = (2{,}0\ \mu\text{F})(24{,}0\ \text{V}) = 48\ [\mu\text{C}]
\]

\textbf{Condensadores en serie con la fuente:}
\[
    \boxed{q_1 = q_6 = Q = 48\ [\mu\text{C}]}
\]
\[
    V_1 = \frac{q_1}{C_1} = \frac{48}{12{,}0} = 4{,}0\ [\text{V}],\qquad
    V_6 = \frac{q_6}{C_6} = \frac{48}{4{,}0} = 12{,}0\ [\text{V}],\qquad
    V_{bd} = \frac{Q}{C_A} = \frac{48}{6{,}0} = 8{,}0\ [\text{V}]
\]
\noindent\textit{Verificaci\'on (serie):} $V_1 + V_{bd} + V_6 = 4{,}0 + 8{,}0 + 12{,}0 = 24{,}0\ [\text{V}] = V$ \checkmark

\medskip
\textbf{Grupo $b$--$d$.} $C_3$ y el camino $C_a$ est\'an sometidos a $V_{bd} = 8{,}0\ [\text{V}]$:
\[
    \boxed{q_3 = C_3\,V_{bd} = (4{,}5)(8{,}0) = 36\ [\mu\text{C}]}\qquad (V_3 = V_{bd} = 8{,}0\ [\text{V}])
\]
\[
    q_a = C_a\,V_{bd} = (1{,}5)(8{,}0) = 12\ [\mu\text{C}]
\]
Como $C_2$, $C_4$ y $C_5$ est\'an en serie, los tres llevan la carga $q_a$:
\[
    \boxed{q_2 = q_4 = q_5 = 12\ [\mu\text{C}]}
\]
\[
    V_2 = \frac{q_2}{C_2} = \frac{12}{6{,}0} = 2{,}0\ [\text{V}],\qquad
    V_4 = \frac{q_4}{C_4} = \frac{12}{3{,}0} = 4{,}0\ [\text{V}],\qquad
    V_5 = \frac{q_5}{C_5} = \frac{12}{6{,}0} = 2{,}0\ [\text{V}]
\]

\noindent\textbf{Verificaci\'on:}
\begin{itemize}\setlength\itemsep{1pt}
    \item Ramas en paralelo: $q_3 + q_a = 36 + 12 = 48\ [\mu\text{C}] = Q$ \checkmark
    \item Camino en serie: $V_2 + V_4 + V_5 = 2{,}0 + 4{,}0 + 2{,}0 = 8{,}0\ [\text{V}] = V_{bd}$ \checkmark
\end{itemize}

\subsection*{Resumen de resultados}

\begin{center}
\renewcommand{\arraystretch}{1.2}
\begin{tabular}{|c|c|c|c|}
\hline
\textbf{Condensador} & \textbf{$C$ [$\mu$F]} & \textbf{Carga $q$ [$\mu$C]} & \textbf{Voltaje $V$ [V]} \\
\hline
$C_1$ & 12{,}0 & 48 & 4{,}0 \\
$C_2$ & 6{,}0  & 12 & 2{,}0 \\
$C_3$ & 4{,}5  & 36 & 8{,}0 \\
$C_4$ & 3{,}0  & 12 & 4{,}0 \\
$C_5$ & 6{,}0  & 12 & 2{,}0 \\
$C_6$ & 4{,}0  & 48 & 12{,}0 \\
\hline
$C_{eq}$ & 2{,}0 & $Q = 48$ & $V = 24{,}0$ \\
\hline
\end{tabular}
\end{center}

\newpage

%========================================================
% PROBLEMA 2  (picofaradios, 3 mallas)
%========================================================
\section*{Problema 2 -- Circuito mixto de siete condensadores}

\subsection*{Enunciado}

En el circuito de la Fig.\ 2, una fuente $V = 18{,}0\ [\text{V}]$ alimenta siete condensadores
con capacitancias $C_1 = 45\ [\text{pF}]$, $C_2 = 30\ [\text{pF}]$, $C_3 = 40\ [\text{pF}]$,
$C_4 = 60\ [\text{pF}]$, $C_5 = 60\ [\text{pF}]$, $C_6 = 40\ [\text{pF}]$ y
$C_7 = 120\ [\text{pF}]$. Determine:

\begin{enumerate}[(a)]
    \item La capacitancia equivalente del circuito.
    \item La energ\'ia almacenada en cada condensador.
\end{enumerate}

\begin{center}
\begin{circuitikz}[every node/.style={font=\small}]
  % fuente (+ arriba) entre a y f
  \draw (0,5) to[battery1, l_=$V$] (0,0);
  \bplus{0.32,2.75}
  % riel superior: a -- C1 -- b -- C2 -- c
  \draw (0,5) to[C, l=$C_1$] (4,5) to[C, l=$C_2$] (8,5);
  % ramas superiores: b -- C3 -- d  y  c -- C4 -- e
  \draw (4,5) to[C, l=$C_3$] (4,2.5);
  \draw (8,5) to[C, l=$C_4$] (8,2.5);
  % cable central d -- e
  \draw (4,2.5) -- (8,2.5);
  % ramas inferiores: d -- C5 -- h  y  e -- C6 -- g
  \draw (4,2.5) to[C, l=$C_5$] (4,0);
  \draw (8,2.5) to[C, l=$C_6$] (8,0);
  % riel inferior: f -- h -- C7 -- g
  \draw (0,0) -- (4,0) to[C, l_=$C_7$] (8,0);
  % nodos
  \node[circ] at (4,5) {}; \node[circ] at (4,2.5) {};
  \node[circ] at (8,2.5) {}; \node[circ] at (4,0) {};
  \node[above left]  at (0,5)   {$a$};
  \node[above right] at (4,5)   {$b$};
  \node[above right] at (8,5)   {$c$};
  \node[below left]  at (4,2.5) {$d$};
  \node[right]       at (8.1,2.5) {$e$};
  \node[below left]  at (0,0)   {$f$};
  \node[below]       at (4,-0.05) {$h$};
  \node[below right] at (8,0)   {$g$};
\end{circuitikz}\\[2pt]
{\footnotesize Fig.\ 2. Circuito mixto de siete condensadores (tres mallas).}
\end{center}

\subsection*{Datos}

\[
    C_1 = 45\ [\text{pF}],\quad C_2 = 30\ [\text{pF}],\quad C_3 = 40\ [\text{pF}],\quad
    C_4 = 60\ [\text{pF}],
\]
\[
    C_5 = 60\ [\text{pF}],\quad C_6 = 40\ [\text{pF}],\quad C_7 = 120\ [\text{pF}],\qquad
    V = 18{,}0\ [\text{V}]
\]

\subsection*{An\'alisis del circuito}

El \textbf{cable} que une $d$ con $e$ hace que ambos puntos est\'en al mismo potencial: forman un
\'unico nodo $M = \{d,\,e\}$. Adem\'as, $c$ y $g$ son esquinas (s\'olo unen dos condensadores) y
$h$ est\'a unido por cable con $f$. Por lo tanto:
\begin{itemize}\setlength\itemsep{1pt}
    \item \textbf{Entre $b$ y $M$} hay dos caminos: $C_3$ directo ($b\to d$) y $C_2$ en serie con
          $C_4$ ($b\to c\to e$). Es decir, $C_3 \parallel (C_2 - C_4)$.
    \item \textbf{Entre $M$ y $h$} hay dos caminos: $C_5$ directo ($d\to h$) y $C_6$ en serie con
          $C_7$ ($e\to g\to h$). Es decir, $C_5 \parallel (C_6 - C_7)$.
    \item $C_1$ ($a$--$b$), el grupo $b$--$M$ y el grupo $M$--$h$ quedan en \textbf{serie} con la
          fuente.
\end{itemize}
\[
    C_{eq}:\qquad C_1 \;-\; \underbrace{\bigl[\,C_3 \parallel (C_2 - C_4)\,\bigr]}_{\text{entre } b \text{ y } M}
    \;-\; \underbrace{\bigl[\,C_5 \parallel (C_6 - C_7)\,\bigr]}_{\text{entre } M \text{ y } h}
\]

\subsection*{Parte (a) -- Capacitancia equivalente}

\textbf{Paso 1.} Pares en serie:
\begin{align*}
    C_a &= \frac{C_2\,C_4}{C_2 + C_4} = \frac{(30)(60)}{30 + 60} = 20\ [\text{pF}]
          \qquad(C_2 - C_4)\\[4pt]
    C_b &= \frac{C_6\,C_7}{C_6 + C_7} = \frac{(40)(120)}{40 + 120} = 30\ [\text{pF}]
          \qquad(C_6 - C_7)
\end{align*}

\textbf{Paso 2.} Grupos en paralelo:
\begin{align*}
    C_A &= C_3 + C_a = 40 + 20 = 60\ [\text{pF}] \qquad(\text{entre } b \text{ y } M)\\[2pt]
    C_B &= C_5 + C_b = 60 + 30 = 90\ [\text{pF}] \qquad(\text{entre } M \text{ y } h)
\end{align*}

\textbf{Paso 3.} $C_1$, $C_A$ y $C_B$ en serie:
\begin{align*}
    \frac{1}{C_{eq}} &= \frac{1}{C_1} + \frac{1}{C_A} + \frac{1}{C_B}
                      = \frac{1}{45} + \frac{1}{60} + \frac{1}{90}\ [\text{pF}^{-1}]\\[4pt]
                     &= \frac{4 + 3 + 2}{180} = \frac{1}{20}\ [\text{pF}^{-1}]
\end{align*}
\[
    \boxed{C_{eq} = 20\ [\text{pF}]}
\]

\begin{center}
\begin{circuitikz}[every node/.style={font=\small}]
  \draw (0,2.4) to[battery1, l_=$V$] (0,0);
  \bplus{0.32,1.45}
  \draw (0,2.4) to[C, l=$C_1$] (3,2.4) to[C, l=$C_A$] (6,2.4)
        to[C, l=$C_B$] (6,0) -- (0,0);
\end{circuitikz}\\[2pt]
{\footnotesize Circuito equivalente en serie: $C_1$, $C_A = C_3 \parallel C_a$ y $C_B = C_5 \parallel C_b$.}
\end{center}

\subsection*{Parte (b) -- Energ\'ia almacenada en cada condensador}

\noindent\textbf{Paso 1: carga y voltaje de cada condensador.}

\textbf{Carga total} (com\'un a $C_1$, $C_A$ y $C_B$, que est\'an en serie):
\[
    Q = C_{eq}\,V = (20\ \text{pF})(18{,}0\ \text{V}) = 360\ [\text{pC}]
    \qquad\Rightarrow\qquad \boxed{q_1 = Q = 360\ [\text{pC}]}
\]
\[
    V_1 = \frac{Q}{C_1} = \frac{360}{45} = 8{,}0\ [\text{V}],\qquad
    V_{bM} = \frac{Q}{C_A} = \frac{360}{60} = 6{,}0\ [\text{V}],\qquad
    V_{Mh} = \frac{Q}{C_B} = \frac{360}{90} = 4{,}0\ [\text{V}]
\]
\noindent\textit{Verificaci\'on (serie):} $V_1 + V_{bM} + V_{Mh} = 8{,}0 + 6{,}0 + 4{,}0 = 18{,}0\ [\text{V}] = V$ \checkmark

\medskip
\textbf{Grupo $b$--$M$} (voltaje $V_{bM} = 6{,}0\ [\text{V}]$):
\[
    \boxed{q_3 = C_3\,V_{bM} = (40)(6{,}0) = 240\ [\text{pC}]}
    \qquad
    \boxed{q_2 = q_4 = C_a\,V_{bM} = (20)(6{,}0) = 120\ [\text{pC}]}
\]
con $V_2 = q_2/C_2 = 4{,}0\ [\text{V}]$ y $V_4 = q_4/C_4 = 2{,}0\ [\text{V}]$.

\medskip
\textbf{Grupo $M$--$h$} (voltaje $V_{Mh} = 4{,}0\ [\text{V}]$):
\[
    \boxed{q_5 = C_5\,V_{Mh} = (60)(4{,}0) = 240\ [\text{pC}]}
    \qquad
    \boxed{q_6 = q_7 = C_b\,V_{Mh} = (30)(4{,}0) = 120\ [\text{pC}]}
\]
con $V_6 = q_6/C_6 = 3{,}0\ [\text{V}]$ y $V_7 = q_7/C_7 = 1{,}0\ [\text{V}]$.

\medskip
\noindent\textbf{Verificaci\'on:}
\begin{itemize}\setlength\itemsep{1pt}
    \item Grupo $b$--$M$: $q_3 + q_a = 240 + 120 = 360\ [\text{pC}] = Q$ \ y \
          $V_2 + V_4 = 4{,}0 + 2{,}0 = 6{,}0\ [\text{V}] = V_{bM}$ \checkmark
    \item Grupo $M$--$h$: $q_5 + q_b = 240 + 120 = 360\ [\text{pC}] = Q$ \ y \
          $V_6 + V_7 = 3{,}0 + 1{,}0 = 4{,}0\ [\text{V}] = V_{Mh}$ \checkmark
\end{itemize}

\medskip
\noindent\textbf{Paso 2: energ\'ia de cada condensador.}

La energ\'ia de cada condensador es $U = \tfrac12\,q\,V$ (con $q$ en pC y $V$ en V, $U$ resulta en pJ):
\begin{align*}
    U_1 &= \tfrac12\,(360)(8{,}0) = 1440\ [\text{pJ}] &
    U_2 &= \tfrac12\,(120)(4{,}0) = 240\ [\text{pJ}]\\[2pt]
    U_3 &= \tfrac12\,(240)(6{,}0) = 720\ [\text{pJ}] &
    U_4 &= \tfrac12\,(120)(2{,}0) = 120\ [\text{pJ}]\\[2pt]
    U_5 &= \tfrac12\,(240)(4{,}0) = 480\ [\text{pJ}] &
    U_6 &= \tfrac12\,(120)(3{,}0) = 180\ [\text{pJ}]\\[2pt]
    U_7 &= \tfrac12\,(120)(1{,}0) = 60\ [\text{pJ}] & &
\end{align*}
\[
    \boxed{U_{\text{total}} = \sum_{i=1}^{7} U_i = 3240\ [\text{pJ}] = 3{,}24\ [\text{nJ}]}
\]

\noindent\textit{Verificaci\'on:} la energ\'ia total tambi\'en se obtiene del circuito equivalente,
\[
    U_{\text{total}} = \tfrac12\,C_{eq}\,V^2 = \tfrac12\,(20\times10^{-12})(18{,}0)^2
                     = 3{,}24\times10^{-9}\ [\text{J}] = 3{,}24\ [\text{nJ}]\ \checkmark
\]
El condensador que m\'as energ\'ia almacena es $C_1$: lleva toda la carga $Q$ y soporta el mayor
voltaje.

\subsection*{Resumen de resultados}

\begin{center}
\renewcommand{\arraystretch}{1.2}
\begin{tabular}{|c|c|c|c|c|}
\hline
\textbf{Condensador} & \textbf{$C$ [pF]} & \textbf{Carga $q$ [pC]} & \textbf{Voltaje $V$ [V]} & \textbf{Energ\'ia $U$ [pJ]} \\
\hline
$C_1$ & 45  & 360 & 8{,}0 & 1440 \\
$C_2$ & 30  & 120 & 4{,}0 & 240 \\
$C_3$ & 40  & 240 & 6{,}0 & 720 \\
$C_4$ & 60  & 120 & 2{,}0 & 120 \\
$C_5$ & 60  & 240 & 4{,}0 & 480 \\
$C_6$ & 40  & 120 & 3{,}0 & 180 \\
$C_7$ & 120 & 120 & 1{,}0 & 60 \\
\hline
$C_{eq}$ & 20 & $Q = 360$ & $V = 18{,}0$ & $U_{\text{total}} = 3240$ \\
\hline
\end{tabular}
\end{center}


\newpage

%========================================================
% PROBLEMA 3  (Guia 8, Problema 8)
%========================================================
\section*{Problema 3 -- Gu\'ia 8, Problema 8: corriente en $R_7$}

\subsection*{Enunciado}

En el circuito de la figura las fuentes son ideales, con $\varepsilon_1 = 15\ [\text{V}]$,
$\varepsilon_2 = 10\ [\text{V}]$ y $\varepsilon_3 = 9\ [\text{V}]$, y
$R_1 = R_2 = R_3 = R_4 = 3\ [\Omega]$, $R_5 = 1\ [\Omega]$, $R_6 = 8\ [\Omega]$ y
$R_7 = 2\ [\Omega]$. Use las leyes de Kirchhoff para determinar la corriente $I$ en $R_7$ (sentido de
referencia indicado por la flecha) e interprete su signo.

\begin{center}
\begin{circuitikz}[every node/.style={font=\small}]
  \ctikzset{bipoles/length=1.15cm}
  % superior: eps1 (+ a la izquierda) y R1
  \draw (0,5) to[battery1, l={$\varepsilon_1 = 15\,\text{V}$}] (3.5,5)
              to[R, l=$R_1$, a=$3\,\Omega$] (7,5);
  \bplus{1.5,4.68}
  % izquierda: R4 y R2; abajo eps3 (+ arriba)
  \draw (0,1.7) to[R, l=$R_4$, a=$3\,\Omega$] (0,3.35)
                to[R, l=$R_2$, a=$3\,\Omega$] (0,5);
  \draw (0,1.7) to[battery1, l_={$\varepsilon_3 = 9\,\text{V}$}] (0,0);
  \bplus{0.32,1.1}
  % central: eps2 (+ arriba) y R5
  \draw (3.5,5) to[battery1, l_={$\varepsilon_2 = 10\,\text{V}$}] (3.5,3.5) -- (3.5,1.7);
  \bplus{3.82,4.5}
  \draw (3.5,0) to[R, l=$R_5$, a=$1\,\Omega$] (3.5,1.7);
  % derecha: R3
  \draw (7,0) -- (7,3.3) to[R, l=$R_3$, a=$3\,\Omega$] (7,5);
  % inferior: R6 y R7 (flecha I de referencia)
  \draw (0,0) to[R, l=$R_6$, a=$8\,\Omega$] (3.5,0)
              to[R, l={$R_7 = 2\,\Omega$}, f>_={$I$}] (7,0);
  % nodos
  \node[circ] at (3.5,5) {}; \node[above right] at (3.5,5.05) {$A$};
  \node[circ] at (3.5,0) {}; \node[below left] at (3.45,-0.05) {$B$};
  % corrientes de rama (rojo)
  \draw[->,red,line width=1.1pt] (0.9,5.32)--(0.3,5.32) node[midway,above,red]{$I_1$};
  \draw[->,red,line width=1.1pt] (3.85,2.2)--(3.85,3.0) node[midway,right,red]{$I_2$};
  \draw[->,red,line width=1.1pt] (7.35,1.2)--(7.35,2.0) node[midway,right,red]{$I_3$};
  % mallas
  \mallah{1.9,2.5}{$M_1$}
  \mallah{5.25,2.5}{$M_2$}
\end{circuitikz}\\[2pt]
{\footnotesize Gu\'ia 8, Fig.\ 8, con corrientes de rama (rojo) y mallas en sentido horario (azul).}
\end{center}

\subsection*{An\'alisis del circuito}

Hay dos nodos, $A$ (arriba) y $B$ (abajo), unidos por tres ramas:
\begin{itemize}\setlength\itemsep{1pt}
    \item \textbf{Rama izquierda} ($\varepsilon_1$, $R_2$, $R_4$, $\varepsilon_3$, $R_6$): corriente
          $I_1$, que sale de $A$ por $\varepsilon_1$ y llega a $B$ por $R_6$.
    \item \textbf{Rama central} ($R_5$, $\varepsilon_2$): corriente $I_2$, que sube de $B$ a $A$.
    \item \textbf{Rama derecha} ($R_7$, $R_3$, $R_1$): corriente $I_3$ en el sentido de la flecha
          ($B\to R_7\to R_3\to R_1\to A$). As\'i, $I = I_3$.
\end{itemize}

\noindent\textit{Convenci\'on para KVL:} en una resistencia recorrida en el sentido de su corriente el
potencial baja ($-IR$), y en contra sube ($+IR$); al atravesar una fuente de $-$ a $+$ el potencial
sube ($+\varepsilon$), y de $+$ a $-$ baja ($-\varepsilon$).

\subsection*{Ecuaciones de Kirchhoff}

\textbf{KCL --- Nodo $A$:} $I_2$ e $I_3$ llegan a $A$ y sale $I_1$:
\[
    \boxed{I_1 = I_2 + I_3} \tag{N}
\]

\textbf{KVL --- Malla $M_1$}, recorrida $B\to R_6\to\varepsilon_3\to R_4\to R_2\to\varepsilon_1\to A\to\varepsilon_2\to R_5\to B$.
En $R_6$, $R_4$, $R_2$ y $R_5$ se avanza contra la corriente; $\varepsilon_3$ se atraviesa de $-$ a
$+$, y $\varepsilon_1$ y $\varepsilon_2$ de $+$ a $-$:
\begin{align*}
    I_1R_6 + \varepsilon_3 + I_1R_4 + I_1R_2 - \varepsilon_1 - \varepsilon_2 + I_2R_5 &= 0\\[2pt]
    I_1\,(R_2+R_4+R_6) + I_2\,R_5 &= \varepsilon_1 + \varepsilon_2 - \varepsilon_3
\end{align*}
\[
    \boxed{14\,I_1 + I_2 = 16} \tag{M1}
\]

\textbf{KVL --- Malla $M_2$}, recorrida $A\to R_1\to R_3\to R_7\to B\to R_5\to\varepsilon_2\to A$.
En $R_1$, $R_3$ y $R_7$ se avanza contra $I_3$; en $R_5$ con $I_2$; $\varepsilon_2$ se atraviesa
de $-$ a $+$:
\begin{align*}
    I_3R_1 + I_3R_3 + I_3R_7 - I_2R_5 + \varepsilon_2 &= 0\\[2pt]
    I_3\,(R_1+R_3+R_7) - I_2\,R_5 &= -\varepsilon_2
\end{align*}
\[
    \boxed{8\,I_3 - I_2 = -10} \tag{M2}
\]

\subsection*{Resoluci\'on del sistema}

Con $I_2 = I_1 - I_3$ (de \textbf{N}) en \textbf{(M1)} y \textbf{(M2)}:
\begin{align*}
    15\,I_1 - I_3 &= 16\\[2pt]
    -I_1 + 9\,I_3 &= -10 \;\;\Rightarrow\;\; I_1 = 9\,I_3 + 10
\end{align*}
Reemplazando: $15\,(9\,I_3 + 10) - I_3 = 16 \;\Rightarrow\; 134\,I_3 = -134$:
\[
    I_3 = -1{,}0\ [\text{A}],\qquad I_1 = 1{,}0\ [\text{A}],\qquad I_2 = I_1 - I_3 = 2{,}0\ [\text{A}]
\]
\[
    \boxed{I = I_3 = -1{,}0\ [\text{A}]}
\]
El signo \textbf{negativo} indica que la corriente en $R_7$ circula en \textbf{sentido contrario} a la
flecha: $1{,}0\ [\text{A}]$ desde la esquina derecha hacia $B$.

\subsection*{Verificaci\'on (balance de potencia)}

$I_1$ atraviesa $\varepsilon_1$ de $-$ a $+$ (entrega) y $\varepsilon_3$ de $+$ a $-$ (absorbe);
$I_2$ atraviesa $\varepsilon_2$ de $-$ a $+$ (entrega):
\begin{align*}
    P_{\text{fuentes}} &= \varepsilon_1 I_1 + \varepsilon_2 I_2 - \varepsilon_3 I_1 = 15 + 20 - 9 = 26\ [\text{W}]\\[2pt]
    P_{\text{disipada}} &= I_1^2(R_2+R_4+R_6) + I_2^2R_5 + I_3^2(R_1+R_3+R_7) = 14 + 4 + 8 = 26\ [\text{W}]\ \checkmark
\end{align*}

\newpage

%========================================================
% PROBLEMA 4  (Guia 8, Problema 10)
%========================================================
\section*{Problema 4 -- Gu\'ia 8, Problema 10: potencia total disipada}

\subsection*{Enunciado}

En el circuito de la figura las fuentes son ideales, con $\varepsilon_1 = 18\ [\text{V}]$,
$\varepsilon_2 = 6\ [\text{V}]$ y $\varepsilon_3 = 2\ [\text{V}]$, y
$R_1 = 2$, $R_2 = 5$, $R_3 = 8$, $R_4 = 4$, $R_5 = 10$, $R_6 = 3$, $R_7 = 12$ y $R_8 = 4\ [\Omega]$.
Calcule la potencia total disipada en todas las resistencias del circuito.

\begin{center}
\begin{circuitikz}[every node/.style={font=\small}]
  \ctikzset{bipoles/length=1.15cm}
  % superior: R1, R2 y R3
  \draw (0,6.2) -- (0.5,6.2) to[R, l=$R_1$, a=$2\,\Omega$] (2.3,6.2)
                to[R, l=$R_2$, a=$5\,\Omega$] (4.2,6.2)
                to[R, l=$R_3$, a=$8\,\Omega$] (6.5,6.2);
  % izquierda superior: R6 abajo, R4 arriba
  \draw (0,3.1) to[R, l=$R_6$, a=$3\,\Omega$] (0,4.65)
                to[R, l=$R_4$, a=$4\,\Omega$] (0,6.2);
  % izquierda inferior: R7 y eps3 (+ arriba)
  \draw (0,1.55) to[R, l=$R_7$, a=$12\,\Omega$] (0,3.1);
  \draw (0,1.55) to[battery1, l_={$\varepsilon_3 = 2\,\text{V}$}] (0,0);
  \bplus{0.32,1.0}
  % rama interior: eps1 (+ abajo)
  \draw (0,3.1) -- (3.5,3.1);
  \draw (3.5,0) to[battery1, l_={$\varepsilon_1 = 18\,\text{V}$}] (3.5,3.1);
  \bplus{3.18,1.3}
  % derecha: R5 arriba y eps2 (+ arriba) abajo
  \draw (6.5,3.1) to[R, l=$R_5$, a=$10\,\Omega$] (6.5,6.2);
  \draw (6.5,3.1) to[battery1, l={$\varepsilon_2 = 6\,\text{V}$}] (6.5,0);
  \bplus{6.18,1.8}
  % inferior
  \draw (0,0) to[R, l=$R_8$, a=$4\,\Omega$] (3.5,0) -- (6.5,0);
  % nodos
  \node[circ] at (0,3.1) {}; \node[left] at (-0.08,3.1) {$N$};
  \node[circ] at (3.5,0) {}; \node[below right] at (3.55,-0.05) {$M$};
  % corrientes de rama (rojo)
  \draw[->,red,line width=1.1pt] (6.85,3.3)--(6.85,3.95) node[midway,right,red]{$I_1$};
  \draw[->,red,line width=1.1pt] (1.3,3.42)--(2.3,3.42) node[midway,above,red]{$I_2$};
  \draw[->,red,line width=1.1pt] (1.0,0.35)--(0.35,0.35) node[midway,above,red]{$I_3$};
  % mallas
  \mallah{3.25,4.65}{$M_1$}
  \mallah{1.75,1.7}{$M_2$}
\end{circuitikz}\\[2pt]
{\footnotesize Gu\'ia 8, Fig.\ 10, con corrientes de rama (rojo) y mallas en sentido horario (azul).}
\end{center}

\subsection*{An\'alisis del circuito}

Hay dos nodos, $N$ (izquierda) y $M$ (abajo), unidos por tres ramas:
\begin{itemize}\setlength\itemsep{1pt}
    \item \textbf{Rama superior} ($\varepsilon_2$, $R_5$, $R_3$, $R_2$, $R_1$, $R_4$, $R_6$): corriente
          $I_1$, que sube por $\varepsilon_2$ desde $M$, recorre el borde superior y llega a $N$.
    \item \textbf{Rama interior} ($\varepsilon_1$, sin resistencia): corriente $I_2$, de $N$ hacia $M$.
    \item \textbf{Rama inferior} ($R_8$, $\varepsilon_3$, $R_7$): corriente $I_3$, que va de $M$ por
          $R_8$, sube por $\varepsilon_3$ y llega a $N$.
\end{itemize}
Para la potencia disipada basta conocer $I_1$ e $I_3$, que son las corrientes que recorren las
resistencias.

\subsection*{Ecuaciones de Kirchhoff}

\textbf{KCL --- Nodo $N$:} $I_1$ e $I_3$ llegan a $N$ y sale $I_2$:
\[
    \boxed{I_2 = I_1 + I_3} \tag{N}
\]

\textbf{KVL --- Malla $M_1$} (rama superior + rama interior), recorrida
$N\to R_6\to R_4\to R_1\to R_2\to R_3\to R_5\to\varepsilon_2\to M\to\varepsilon_1\to N$. En todas las
resistencias se avanza contra $I_1$; $\varepsilon_2$ (al bajar) y $\varepsilon_1$ (al subir) se
atraviesan de $+$ a $-$:
\begin{align*}
    I_1\,(R_6+R_4+R_1+R_2+R_3+R_5) - \varepsilon_2 - \varepsilon_1 &= 0\\[2pt]
    32\,I_1 = 24 \;\;\Rightarrow\;\; I_1 &= 0{,}75\ [\text{A}]
\end{align*}

\textbf{KVL --- Malla $M_2$} (rama interior + rama inferior), recorrida
$N\to\varepsilon_1\to M\to R_8\to\varepsilon_3\to R_7\to N$. $\varepsilon_1$ (al bajar) y
$\varepsilon_3$ (al subir) se atraviesan de $-$ a $+$, y en $R_8$ y $R_7$ se avanza con $I_3$:
\begin{align*}
    \varepsilon_1 - I_3R_8 + \varepsilon_3 - I_3R_7 &= 0\\[2pt]
    16\,I_3 = 20 \;\;\Rightarrow\;\; I_3 &= 1{,}25\ [\text{A}]
\end{align*}
y de \textbf{(N)}: $I_2 = 0{,}75 + 1{,}25 = 2{,}0\ [\text{A}]$.

\subsection*{Potencia total disipada}

\begin{align*}
    P_{\text{tot}} &= I_1^2\,(R_1+R_2+R_3+R_4+R_5+R_6) + I_3^2\,(R_7+R_8)\\[2pt]
                   &= (0{,}75)^2(32) + (1{,}25)^2(16) = 18 + 25
\end{align*}
\[
    \boxed{P_{\text{tot}} = 43\ [\text{W}]}
\]

\noindent\textbf{Verificaci\'on:} las tres fuentes entregan energ\'ia, porque cada corriente las
atraviesa de $-$ a $+$:
$\varepsilon_1 I_2 + \varepsilon_2 I_1 + \varepsilon_3 I_3 = 36 + 4{,}5 + 2{,}5 = 43\ [\text{W}]$ \checkmark

\newpage

%========================================================
% PROBLEMA 5  (lazo triangular y carga puntual)
%========================================================
\section*{Problema 5 -- Fuerza magn\'etica sobre una carga y un lazo triangular}

\subsection*{Enunciado}

En la Fig.~3 se muestra un cable en forma de tri\'angulo rect\'angulo cuya longitud del lado $ab$ es
$l_{ab} = 15\ [\text{cm}]$, por donde circula una corriente $I = 2{,}5\ [\text{A}]$ en sentido
horario. En alg\'un instante de tiempo, una carga puntual $q = -4{,}5\ [\text{nC}]$ que tiene una
rapidez $v = 300\ [\text{m/s}]$, dirigida en un \'angulo de $20^\circ$ con respecto a la horizontal, se
encuentra cerca del cable. Tanto el cable como la carga se encuentran en una regi\'on de campo
magn\'etico uniforme de magnitud $B = 1{,}5\ [\text{T}]$, cuya direcci\'on se indica en la figura.
Determine:
\begin{enumerate}[(a)]
    \item La fuerza magn\'etica $\vec{F}$ que ejerce $\vec{B}$ sobre la carga puntual.
    \item La fuerza magn\'etica $\vec{F}$ que ejerce $\vec{B}$ sobre cada lado del tri\'angulo $abc$.
    \item El momento magn\'etico $\vec{\mu}$ del lazo de corriente.
    \item El torque $\vec{\tau}$ que ejerce $\vec{B}$ sobre el lazo de corriente.
\end{enumerate}

\begin{center}
\begin{tikzpicture}[>=Stealth, line width=0.8pt, scale=1.25]
    % lineas de campo a 25 grados
    \foreach \x in {-1.5,0.3,2.1,3.9,5.7}
        \draw[->,blue!65] (\x,-0.4) -- ++({4.2*cos(25)},{4.2*sin(25)});
    \node[blue!60!black] at (5.6,3.0) {$\vec{B}$};
    \draw[dashed,gray] (2.7,1.4) -- (3.7,1.4);
    \draw[gray] ({2.7+0.65},1.4) arc[start angle=0, end angle=25, radius=0.65];
    \node[font=\footnotesize] at (3.65,1.65) {$25^\circ$};
    % triangulo (a recto, angulo 15 en b)
    \coordinate (a) at (0,0); \coordinate (b) at (4.5,0); \coordinate (c) at (0,{4.5*tan(15)});
    \draw[line width=1.5pt] (a) -- (b) -- (c) -- cycle;
    \node[below left] at (a) {$a$}; \node[below right] at (b) {$b$}; \node[above left] at (c) {$c$};
    \draw[gray] ($(b)+(180:0.9)$) arc[start angle=180, end angle=165, radius=0.9];
    \node[font=\footnotesize] at ($(b)+(172:1.25)$) {$15^\circ$};
    % corriente horaria c->b->a->c
    \draw[->,red,line width=1.2pt] ($(c)!0.3!(b)+(0,0.18)$) -- ($(c)!0.6!(b)+(0,0.18)$);
    \draw[->,red,line width=1.2pt] (3.0,-0.22) -- (1.6,-0.22);
    \draw[->,red,line width=1.2pt] (-0.22,0.25) -- (-0.22,0.95);
    \node[red] at (2.2,-0.55) {$I$};
    % carga y velocidad (20 grados sobre -x)
    \fill[orange!50] (6.6,1.3) circle (0.22); \draw (6.6,1.3) circle (0.22); \node at (6.6,1.3) {\scriptsize $q$};
    \draw[->,line width=1.2pt] (6.6,1.3) -- ++({1.5*cos(160)},{1.5*sin(160)}) node[above]{$\vec{v}$};
    \draw[dashed,gray] (6.6,1.3) -- (5.3,1.3);
    \draw[gray] ({6.6-0.75},1.3) arc[start angle=180, end angle=160, radius=0.75];
    \node[font=\footnotesize] at (5.55,1.12) {$20^\circ$};
    % ejes
    \begin{scope}[shift={(-1.8,2.2)}]
        \draw[->] (0,0) -- (0.8,0) node[right]{$x$};
        \draw[->] (0,0) -- (0,0.8) node[above]{$y$};
        \draw (0,0) circle (0.1); \fill (0,0) circle (0.035); \node[below left] at (0,0) {$z$};
    \end{scope}
\end{tikzpicture}\\[2pt]
{\footnotesize Fig.\ 3 (redibujada; $\hat{k}$ sale de la p\'agina).}
\end{center}

\noindent\textbf{Vectores del campo} (de la figura):
\[
    \vec{B} = B\left(\cos25^\circ\,\hat{\imath} + \sin25^\circ\,\hat{\jmath}\right)
            = 1{,}359\,\hat{\imath} + 0{,}634\,\hat{\jmath}\ [\text{T}]
\]

\subsection*{(a) Fuerza sobre la carga puntual}

La velocidad forma $20^\circ$ sobre el eje $-x$:
$\vec{v} = v\left(-\cos20^\circ\,\hat{\imath} + \sin20^\circ\,\hat{\jmath}\right)$. Producto cruz t\'ermino
a t\'ermino:
\begin{align*}
    \vec{v}\times\vec{B} &= vB\bigl[-\cos20^\circ\cos25^\circ(\hat{\imath}\times\hat{\imath})
                            - \cos20^\circ\sin25^\circ(\hat{\imath}\times\hat{\jmath})\\
                         &\qquad\;\; + \sin20^\circ\cos25^\circ(\hat{\jmath}\times\hat{\imath})
                            + \sin20^\circ\sin25^\circ(\hat{\jmath}\times\hat{\jmath})\bigr]\\
                         &= -vB\left(\cos20^\circ\sin25^\circ + \sin20^\circ\cos25^\circ\right)\hat{k}
                          = -vB\sin45^\circ\,\hat{k}\\
                         &= -(300)(1{,}5)(0{,}7071)\,\hat{k} = -318{,}2\,\hat{k}\ [\text{T m/s}]
\end{align*}
\[
    \vec{F} = q\,\vec{v}\times\vec{B} = (-4{,}5\times10^{-9})(-318{,}2)\,\hat{k}
    \qquad\Rightarrow\qquad
    \boxed{\vec{F} = 1{,}43\times10^{-6}\,\hat{k}\ [\text{N}]}
\]
La fuerza sale de la p\'agina (la carga es negativa).

\subsection*{(b) Fuerza sobre cada lado del tri\'angulo}

\textbf{Geometr\'ia} (\'angulo recto en $a$, $15^\circ$ en $b$):
\[
    l_{ac} = l_{ab}\tan15^\circ = 0{,}0402\ [\text{m}],\qquad
    l_{bc} = \frac{l_{ab}}{\cos15^\circ} = 0{,}1553\ [\text{m}]
\]
Con $a=(0,0)$, $b=(0{,}15;\,0)$ y $c=(0;\,0{,}0402)$ y la corriente horaria $c\to b\to a\to c$:
\[
    \vec{l}_{cb} = 0{,}15\,\hat{\imath} - 0{,}0402\,\hat{\jmath},\qquad
    \vec{l}_{ba} = -0{,}15\,\hat{\imath},\qquad
    \vec{l}_{ac} = 0{,}0402\,\hat{\jmath}\ [\text{m}]
\]
Fuerzas $\vec{F} = I\,\vec{l}\times\vec{B}$:
\begin{align*}
    \vec{F}_{cb} &= 2{,}5\bigl[(0{,}15)(0{,}634)(\hat{\imath}\times\hat{\jmath})
                   - (0{,}0402)(1{,}359)(\hat{\jmath}\times\hat{\imath})\bigr]
                  = 2{,}5\,(0{,}0951 + 0{,}0546)\,\hat{k}\\[2pt]
    \vec{F}_{ba} &= 2{,}5\,(-0{,}15)(0{,}634)(\hat{\imath}\times\hat{\jmath})
                  = -2{,}5\,(0{,}0951)\,\hat{k}\\[2pt]
    \vec{F}_{ac} &= 2{,}5\,(0{,}0402)(1{,}359)(\hat{\jmath}\times\hat{\imath})
                  = -2{,}5\,(0{,}0546)\,\hat{k}
\end{align*}
(los t\'erminos $\hat{\imath}\times\hat{\imath}$ y $\hat{\jmath}\times\hat{\jmath}$ son nulos).
\[
    \boxed{\vec{F}_{cb} = 0{,}374\,\hat{k}\ [\text{N}],\qquad
           \vec{F}_{ba} = -0{,}238\,\hat{k}\ [\text{N}],\qquad
           \vec{F}_{ac} = -0{,}137\,\hat{k}\ [\text{N}]}
\]
\noindent\textit{Verificaci\'on:} $\vec{F}_{cb} + \vec{F}_{ba} + \vec{F}_{ac} = (0{,}374 - 0{,}238 - 0{,}137)\,\hat{k}
\approx \vec{0}$ \checkmark\ (lazo cerrado en campo uniforme).

\subsection*{(c) Momento magn\'etico}

La corriente circula en sentido horario: por la regla de la mano derecha $\hat{n} = -\hat{k}$. El
\'area del tri\'angulo es $A = \tfrac12\,l_{ab}\,l_{ac} = \tfrac12(0{,}15)(0{,}0402) = 3{,}01\times10^{-3}\ [\text{m}^2]$:
\[
    \vec{\mu} = I\,A\,\hat{n} = (2{,}5)(3{,}01\times10^{-3})(-\hat{k})
    \qquad\Rightarrow\qquad
    \boxed{\vec{\mu} = -7{,}54\times10^{-3}\,\hat{k}\ [\text{A m}^2]}
\]

\subsection*{(d) Torque sobre el lazo}

\begin{align*}
    \vec{\tau} = \vec{\mu}\times\vec{B}
      &= -7{,}54\times10^{-3}\,\hat{k}\times\left(1{,}359\,\hat{\imath} + 0{,}634\,\hat{\jmath}\right)\\[2pt]
      &= -7{,}54\times10^{-3}\left[1{,}359\,(\hat{k}\times\hat{\imath}) + 0{,}634\,(\hat{k}\times\hat{\jmath})\right]\\[2pt]
      &= -7{,}54\times10^{-3}\left[1{,}359\,\hat{\jmath} - 0{,}634\,\hat{\imath}\right]
\end{align*}
\[
    \boxed{\vec{\tau} = 4{,}78\times10^{-3}\,\hat{\imath} - 1{,}02\times10^{-2}\,\hat{\jmath}\ [\text{N m}]}
\]
Como $\vec{B}$ est\'a en el plano del lazo y $\vec{\mu}$ es perpendicular a \'el, el torque es m\'aximo:
$|\vec{\tau}| = \mu B = (7{,}54\times10^{-3})(1{,}5) = 1{,}13\times10^{-2}\ [\text{N m}]$.

\end{document}
